% ==============================================================================
% System Heterogeneity in DermaMNIST Classification
% ------------------------------------------------------------------------------
% Drop-in section for the §Results chapter, sibling of
% 10_statistical_heterogeneity.tex.
%
% Reporting style matches the federated-learning literature convention
% (descriptive comparisons of means, visual curve comparisons, bolded
% best results in tables), per Li 2020, McMahan 2017, Wang 2020, Hsu 2019.
%
% Scope: FedAvg vs FedProx only. The system-heterogeneity conditions
% are evaluated on TWO partition designs (engineered balanced_paired and
% IID) to give a cross-partition mechanism control.
%
% Required figures (in figures/):
%   F6_system_het_schedules.pdf       -- per-(round, client) epoch heatmaps
%   F7_system_het_convergence.pdf     -- 2x3 convergence panel (eng + IID)
%
% Required tables (in tables/):
%   T06_system_het_headline.tex       -- engineered C0/C1/C2 summary
%   T07_system_het_per_class.tex      -- engineered C2 per-class
%   T08_straggler_tolerance.tex       -- tolerance ratios
%   T09_asymmetric_decomposition.tex  -- Li 2020 §5.2 replication (Δ_total)
%   T12_system_het_iid.tex            -- IID C0/C1/C2 summary
%   T13_cross_partition.tex           -- engineered vs IID side-by-side
%
% Usage in the main Overleaf file:
%   % ==============================================================================
% System Heterogeneity in DermaMNIST Classification
% ------------------------------------------------------------------------------
% Drop-in section for the §Results chapter, sibling of
% 10_statistical_heterogeneity.tex.
%
% Reporting style matches the federated-learning literature convention
% (descriptive comparisons of means, visual curve comparisons, bolded
% best results in tables), per Li 2020, McMahan 2017, Wang 2020, Hsu 2019.
%
% Scope: FedAvg vs FedProx only. The system-heterogeneity conditions
% are evaluated on TWO partition designs (engineered balanced_paired and
% IID) to give a cross-partition mechanism control.
%
% Required figures (in figures/):
%   F6_system_het_schedules.pdf       -- per-(round, client) epoch heatmaps
%   F7_system_het_convergence.pdf     -- 2x3 convergence panel (eng + IID)
%
% Required tables (in tables/):
%   T06_system_het_headline.tex       -- engineered C0/C1/C2 summary
%   T07_system_het_per_class.tex      -- engineered C2 per-class
%   T08_straggler_tolerance.tex       -- tolerance ratios
%   T09_asymmetric_decomposition.tex  -- Li 2020 §5.2 replication (Δ_total)
%   T12_system_het_iid.tex            -- IID C0/C1/C2 summary
%   T13_cross_partition.tex           -- engineered vs IID side-by-side
%
% Usage in the main Overleaf file:
%   % ==============================================================================
% System Heterogeneity in DermaMNIST Classification
% ------------------------------------------------------------------------------
% Drop-in section for the §Results chapter, sibling of
% 10_statistical_heterogeneity.tex.
%
% Reporting style matches the federated-learning literature convention
% (descriptive comparisons of means, visual curve comparisons, bolded
% best results in tables), per Li 2020, McMahan 2017, Wang 2020, Hsu 2019.
%
% Scope: FedAvg vs FedProx only. The system-heterogeneity conditions
% are evaluated on TWO partition designs (engineered balanced_paired and
% IID) to give a cross-partition mechanism control.
%
% Required figures (in figures/):
%   F6_system_het_schedules.pdf       -- per-(round, client) epoch heatmaps
%   F7_system_het_convergence.pdf     -- 2x3 convergence panel (eng + IID)
%
% Required tables (in tables/):
%   T06_system_het_headline.tex       -- engineered C0/C1/C2 summary
%   T07_system_het_per_class.tex      -- engineered C2 per-class
%   T08_straggler_tolerance.tex       -- tolerance ratios
%   T09_asymmetric_decomposition.tex  -- Li 2020 §5.2 replication (Δ_total)
%   T12_system_het_iid.tex            -- IID C0/C1/C2 summary
%   T13_cross_partition.tex           -- engineered vs IID side-by-side
%
% Usage in the main Overleaf file:
%   % ==============================================================================
% System Heterogeneity in DermaMNIST Classification
% ------------------------------------------------------------------------------
% Drop-in section for the §Results chapter, sibling of
% 10_statistical_heterogeneity.tex.
%
% Reporting style matches the federated-learning literature convention
% (descriptive comparisons of means, visual curve comparisons, bolded
% best results in tables), per Li 2020, McMahan 2017, Wang 2020, Hsu 2019.
%
% Scope: FedAvg vs FedProx only. The system-heterogeneity conditions
% are evaluated on TWO partition designs (engineered balanced_paired and
% IID) to give a cross-partition mechanism control.
%
% Required figures (in figures/):
%   F6_system_het_schedules.pdf       -- per-(round, client) epoch heatmaps
%   F7_system_het_convergence.pdf     -- 2x3 convergence panel (eng + IID)
%
% Required tables (in tables/):
%   T06_system_het_headline.tex       -- engineered C0/C1/C2 summary
%   T07_system_het_per_class.tex      -- engineered C2 per-class
%   T08_straggler_tolerance.tex       -- tolerance ratios
%   T09_asymmetric_decomposition.tex  -- Li 2020 §5.2 replication (Δ_total)
%   T12_system_het_iid.tex            -- IID C0/C1/C2 summary
%   T13_cross_partition.tex           -- engineered vs IID side-by-side
%
% Usage in the main Overleaf file:
%   \input{path/to/11_system_heterogeneity}
% ==============================================================================

\section{System Heterogeneity}
\label{sec:results-system-het}

System heterogeneity refers to differences in the local computational
budget that participating clients contribute per communication round.
In production federated deployments this is the dominant operational
challenge: institutions run on heterogeneous hardware, on-call training
windows differ by site, and not all clients can complete the same
amount of local work within a round's wall-clock deadline. The proximal
term in FedProx was originally motivated for this regime by
\citet{li2020fedprox}: clients may submit $\gamma$-inexact updates ---
partial work that does not fully minimise the local objective --- while
the proximal anchor $(\mu/2)\lVert w - w^t \rVert^2$ provably preserves
convergence by preventing inexact updates from drifting too far from
the broadcast iterate $w^t$ (Theorem~4).

This section reports the FedProx-vs-FedAvg comparison under three
straggler conditions on the engineered \texttt{balanced\_paired\_7\_clients}
partition (the same partition used as the primary in
Section~\ref{sec:results-statistical-het}) and additionally on the IID
partition (\texttt{iid\_7\_clients}) as a mechanism control. We follow
the same descriptive reporting convention as the
statistical-heterogeneity section: mean $\pm$ SD across paired seeds,
visual curve comparisons, bolded best results.

\paragraph{Terminology.}
We use ``system heterogeneity'' as shorthand for
\emph{computational-work heterogeneity}: variability in the number of
local SGD steps each client performs per round. The manipulation acts
only on the algorithm via heterogeneous local-epoch budgets $\tau_i$;
we do not simulate hardware speed, network latency, or device
availability.

\subsection{Experimental design and condition motivation}
\label{sec:system-het-conditions}

Three system-heterogeneity conditions are evaluated. Each schedule is
deterministic given the seed (see \texttt{fl/system\_het.py}), so
paired FedAvg/FedProx runs share an identical per-(round, client)
local-epoch assignment and the within-pair $\Delta$ reflects only the
algorithmic difference. The schedules under seed~42 are visualised in
Figure~\ref{fig:system-het-schedules}.

\paragraph{C0 (no system heterogeneity, baseline).}
Every client performs $E = 20$ local epochs every round. The design
is identical to the headline statistical-heterogeneity setup, but
re-executed in the Flower runtime so that all three system-het
conditions share a single runtime. C0 provides the within-runtime
baseline against which the system-het amplification can be assessed.

\paragraph{C1 (fixed stragglers).}
Two designated clients (C5 and C6) always perform
$E_{\text{straggler}} = 5$ local epochs; the remaining five clients
perform the full $E = 20$. Under the engineered partition, C5 holds
the melanoma/vascular specialty pair and C6 is the nevi-only
generalist, so the C1 design is structurally coupled with the
clinical-role assignment under that partition. Under the IID partition
the choice of (C5, C6) is arbitrary (clients are exchangeable), so the
IID C1 isolates the pure compute-heterogeneity effect more cleanly.

\paragraph{C2 (random stragglers).}
Each communication round, $\lfloor 0.5 \times 7 \rfloor + 1 = 4$
clients are randomly designated stragglers; each straggler performs
$E_i \sim \mathrm{Uniform}\{1, \dots, 19\}$ local epochs and the
remaining three perform $E = 20$. Aggregated across the ten seeds,
each client serves as a straggler in approximately $57\%$ of rounds.
This follows the protocol of \citet{li2020fedprox} §5.2 up to two
implementation details: stragglers always perform at least one
optimiser step (the Li et al.\ protocol allows $E_i = 0$), and both
algorithms include partial updates in aggregation (the Li et al.\
canonical FedAvg drops them; we recover the asymmetric protocol
separately in Section~\ref{sec:asymmetric}).

\begin{figure}[!htbp]
\centering
\includegraphics[width=\textwidth]{F6_system_het_schedules.pdf}
\caption{Per-(round, client) local-epoch schedules under the three
system-heterogeneity conditions, reconstructed deterministically from
seed~42. C0 (left) is uniform $E = 20$. C1 (middle) holds clients
C5 and C6 at $E = 5$ in every round. C2 (right) randomly designates
4 of 7 clients per round as stragglers with
$E_i \sim \mathrm{Uniform}\{1, \dots, 19\}$.}
\label{fig:system-het-schedules}
\end{figure}

\subsection{Results on the engineered partition}
\label{sec:system-het-engineered}

Table~\ref{tab:system-het-headline} reports test macro-F1 under the
three conditions on the engineered partition. The mean within-pair
$\Delta$ rises monotonically from $+0.007$ at the no-system-het
baseline (C0) to $+0.013$ under fixed stragglers (C1) to $+0.017$
under random stragglers (C2). FedProx wins 7--8 of 10 paired seeds
in each condition. The pattern is consistent with the canonical
motivation of \citet{li2020fedprox}: as the amount of inexact local
work increases, the proximal regulariser's drift-control mechanism
has more drift to constrain and the within-pair advantage of FedProx
descriptively grows.

\input{T06_system_het_headline}

\paragraph{Per-class breakdown under C2.}
Table~\ref{tab:system-het-per-class} decomposes the C2 result by
class. The pattern matches that observed in the statistical-
heterogeneity section: the FedProx advantage is concentrated on the
melanoma class ($\Delta = +0.085$ macro-F1, FedProx wins 10/10
paired seeds), while the majority class melanocytic nevi shows
essentially zero difference. The proximal-anchor mechanism is
expected to be active only where inter-client drift exists, and the
melanoma class --- held by two specialist clients under the engineered
partition --- is exactly where drift is structurally produced.

\input{T07_system_het_per_class}

\paragraph{Straggler-tolerance ratios.}
Table~\ref{tab:system-het-tolerance} reports the ratio
$\rho_a^c = M_a^c / M_a^{C0}$ of each algorithm's mean macro-F1 under
the straggler conditions versus its no-system-het baseline. Both
algorithms maintain $\rho \approx 1$ under both straggler conditions,
indicating that the engineered partition's headline macro-F1 is largely
preserved under the introduction of system heterogeneity. FedProx's
$\rho$ is descriptively higher than FedAvg's in both C1 and C2,
consistent with the per-condition $\Delta$ ordering in
Table~\ref{tab:system-het-headline}.

\input{T08_straggler_tolerance}

\begin{figure}[!htbp]
\centering
\includegraphics[width=\textwidth]{F7_system_het_convergence.pdf}
\caption{Validation macro-F1 trajectories under the three
system-heterogeneity conditions on each partition design (Flower
runtime, $n = 10$ paired seeds, mean $\pm$ SEM). \textbf{Top row:}
engineered partition; FedProx (dark slate blue) maintains a visible
upward separation from FedAvg (mint green) throughout training in
all three conditions, with the gap widening from C0 to C2.
\textbf{Bottom row:} IID partition; the two algorithms' trajectories
are essentially superimposed in all three conditions, consistent with
the proximal-anchor mechanism's null prediction under i.i.d.\ data.}
\label{fig:system-het-convergence}
\end{figure}

\subsection{Cross-partition control: IID partition}
\label{sec:system-het-iid}

To verify that the system-heterogeneity amplification observed on the
engineered partition reflects the proximal-anchor mechanism specifically
--- rather than a generic regularisation effect --- we ran the same
three conditions on the IID partition, in which every client receives
a uniform random sample of the global distribution and no inter-client
class skew exists. Under the proximal-anchor mechanism, the algorithm
has no drift to constrain on the IID partition and FedProx is
predicted to provide no benefit, irrespective of straggler condition.

Table~\ref{tab:system-het-iid} reports the IID-partition results. The
within-pair $\Delta$ values are small in absolute magnitude under all
three conditions ($-0.006$ at C0, $+0.010$ at C1, $-0.013$ at C2) and
not consistently signed across the three. The two algorithms' macro-F1
distributions overlap substantially on every condition. Whichever
straggler condition is applied, FedProx and FedAvg perform
indistinguishably on the IID partition --- consistent with the
proximal-anchor mechanism's null prediction.

\input{T12_system_het_iid}

\paragraph{Cross-partition comparison.}
Table~\ref{tab:system-het-cross-partition} contrasts the
within-pair $\Delta$ values between the two partition designs.
On the engineered partition, $\Delta$ is positive in all three
conditions and grows with straggler severity (mechanism active,
amplification visible). On the IID partition, $\Delta$ is small and
not consistently positive (mechanism inert). The pattern is exactly
the cross-partition signature the proximal-anchor mechanism predicts.

\input{T13_cross_partition}

Figure~\ref{fig:system-het-convergence} shows the trajectories
side-by-side: the top row (engineered) consistently displays a
visible FedProx-FedAvg gap, while the bottom row (IID) shows the
two algorithms' trajectories essentially superimposed. The visual
contrast between the two rows is the most direct illustration of
the mechanism's data-dependence: FedProx helps where there is class
skew, and is inert where there is not.

\subsection{Partial client participation}
\label{sec:partial-participation}

The conditions reported above vary the second canonical FL parameter
of system heterogeneity --- the local-epoch budget $E$ --- while
holding the per-round participation fraction fixed at $C = 1.0$ (every
client participates every round). The third canonical FL parameter,
$C$, is itself an axis of system heterogeneity widely tested in the
literature (\citealt{mcmahan2017fedavg} §3.1 varies
$C \in \{0.1, 0.2, 0.5, 1.0\}$; \citealt{li2020fedprox} §5.3 tests
FedProx under partial participation on Synthetic and FEMNIST;
\citealt{kairouz2021advances} lists $C$ as a defining FL setup
parameter). In real federated deployments, not all clients respond
every round, so the full-participation setting reported above is the
cleanest baseline but the least operationally realistic case. We
additionally report a partial-participation condition at $C = 0.5$
(four of the seven clients sampled uniformly at random per round,
computed as $\max(1, \lfloor 0.5 \times 7 \rfloor) = 4$), holding all
other hyperparameters identical to the engineered C0 baseline.

Table~\ref{tab:partial-participation} reports the result. Under
partial participation the within-pair $\Delta$ is
\TODOhpc{[fill in: $X.XXX$ macro-F1, $Y/10$ paired wins once HPC
results land]}, [direction: descriptively positive / comparable to /
descriptively negative versus] the $C = 1.0$ baseline of $+0.007$
macro-F1. The partial-participation setting introduces additional
aggregation-pool variance per round (fewer clients contributing,
varying subsets across rounds), which on mechanism grounds is
expected to interact with the proximal anchor's drift-control
mechanism in a regime not directly tested by the local-epoch
heterogeneity conditions above. The descriptive result here covers
the second canonical FL system-heterogeneity axis and complements the
local-epoch heterogeneity reported in
Sections~\ref{sec:system-het-engineered}--\ref{sec:system-het-iid}.

\input{T14_partial_participation}

\subsection{Replication of the canonical Li et al.\ §5.2 protocol}
\label{sec:asymmetric}

The H1/H2 tests above use a symmetric inclusion policy in which both
FedAvg and FedProx aggregate over all clients including stragglers.
The canonical asymmetric protocol of \citet{li2020fedprox} §5.2, in
contrast, has FedAvg drop straggler updates entirely while FedProx
includes them as $\gamma$-inexact partial work. The original paper
reports a substantial FedProx advantage under this asymmetric setup.

To verify that this canonical setup reproduces qualitatively on
DermaMNIST, we additionally ran the FedAvg-drop arm on the engineered
partition with random stragglers (the C2 schedule). Combined with the
existing FedProx-include arm from the symmetric C2 data on the same
ten seeds, this gives the within-pair comparison in
Table~\ref{tab:asymmetric-replication}.

\input{T09_asymmetric_decomposition}

The FedProx-include configuration outperforms FedAvg-drop by a mean
within-pair difference of $+0.068$ macro-F1, with FedProx winning
all 10 paired seeds. This reproduces the qualitative finding of
\citet{li2020fedprox} §5.2 on the DermaMNIST task: under the
canonical asymmetric straggler protocol, the FedProx configuration
produces a substantially higher macro-F1 than the FedAvg
configuration.

\subsection{Summary}
\label{sec:system-het-summary}

On the engineered partition the within-pair FedProx-vs-FedAvg
difference rises monotonically with straggler severity (C0:
$+0.007$, C1: $+0.013$, C2: $+0.017$), consistent with the canonical
\citet{li2020fedprox} motivation. The C2 per-class breakdown localises
the advantage to the melanoma class ($\Delta = +0.085$ macro-F1, 10/10
paired wins). Under partial client participation ($C = 0.5$,
Section~\ref{sec:partial-participation}) the descriptive result is
\TODOhpc{[fill in once HPC data lands]}. The canonical asymmetric
protocol of \citet{li2020fedprox} §5.2 reproduces qualitatively:
$\Delta_{\text{total}} = +0.068$ macro-F1 with FedProx winning all 10
paired seeds. On the IID partition the within-pair $\Delta$ is small
and not consistently positive under any of the three straggler
conditions, in line with the proximal-anchor mechanism's prediction
that the regulariser is inert when no inter-client class skew exists.

\paragraph{Scope and limitations.}
The within-pair $\Delta$ magnitudes on the engineered partition
($+0.007$--$+0.017$ macro-F1) are smaller than the effect sizes
originally reported in \citet{li2020fedprox} on Synthetic, FEMNIST,
Shakespeare, and Sentiment140 benchmarks. This is consistent with
the small absolute magnitude observed in the statistical-heterogeneity
arm under the same partition and runtime, and is plausibly attributable
to the smaller dataset, more challenging architecture, and
medical-imaging difficulty discussed in
Section~\ref{sec:headline-summary}. The C1 fixed-straggler condition
on the engineered partition is structurally coupled with the clinical
role assignment (C5 holds the melanoma/vascular pair); its descriptive
result should be read with that confound in mind. The asymmetric
protocol of Section~\ref{sec:asymmetric} provides a single point of
comparison against the canonical Li et al.\ setup but is not a primary
analysis; it is included to demonstrate qualitative reproducibility of
their §5.2 result on DermaMNIST. The straggler schedules are
deterministic functions of the seed and have been verified by direct
reconstruction from the JSON provenance --- the empirical per-client
straggler frequency, $E_i$ range, and round counts match the schedule
specifications exactly.

% ==============================================================================

\section{System Heterogeneity}
\label{sec:results-system-het}

System heterogeneity refers to differences in the local computational
budget that participating clients contribute per communication round.
In production federated deployments this is the dominant operational
challenge: institutions run on heterogeneous hardware, on-call training
windows differ by site, and not all clients can complete the same
amount of local work within a round's wall-clock deadline. The proximal
term in FedProx was originally motivated for this regime by
\citet{li2020fedprox}: clients may submit $\gamma$-inexact updates ---
partial work that does not fully minimise the local objective --- while
the proximal anchor $(\mu/2)\lVert w - w^t \rVert^2$ provably preserves
convergence by preventing inexact updates from drifting too far from
the broadcast iterate $w^t$ (Theorem~4).

This section reports the FedProx-vs-FedAvg comparison under three
straggler conditions on the engineered \texttt{balanced\_paired\_7\_clients}
partition (the same partition used as the primary in
Section~\ref{sec:results-statistical-het}) and additionally on the IID
partition (\texttt{iid\_7\_clients}) as a mechanism control. We follow
the same descriptive reporting convention as the
statistical-heterogeneity section: mean $\pm$ SD across paired seeds,
visual curve comparisons, bolded best results.

\paragraph{Terminology.}
We use ``system heterogeneity'' as shorthand for
\emph{computational-work heterogeneity}: variability in the number of
local SGD steps each client performs per round. The manipulation acts
only on the algorithm via heterogeneous local-epoch budgets $\tau_i$;
we do not simulate hardware speed, network latency, or device
availability.

\subsection{Experimental design and condition motivation}
\label{sec:system-het-conditions}

Three system-heterogeneity conditions are evaluated. Each schedule is
deterministic given the seed (see \texttt{fl/system\_het.py}), so
paired FedAvg/FedProx runs share an identical per-(round, client)
local-epoch assignment and the within-pair $\Delta$ reflects only the
algorithmic difference. The schedules under seed~42 are visualised in
Figure~\ref{fig:system-het-schedules}.

\paragraph{C0 (no system heterogeneity, baseline).}
Every client performs $E = 20$ local epochs every round. The design
is identical to the headline statistical-heterogeneity setup, but
re-executed in the Flower runtime so that all three system-het
conditions share a single runtime. C0 provides the within-runtime
baseline against which the system-het amplification can be assessed.

\paragraph{C1 (fixed stragglers).}
Two designated clients (C5 and C6) always perform
$E_{\text{straggler}} = 5$ local epochs; the remaining five clients
perform the full $E = 20$. Under the engineered partition, C5 holds
the melanoma/vascular specialty pair and C6 is the nevi-only
generalist, so the C1 design is structurally coupled with the
clinical-role assignment under that partition. Under the IID partition
the choice of (C5, C6) is arbitrary (clients are exchangeable), so the
IID C1 isolates the pure compute-heterogeneity effect more cleanly.

\paragraph{C2 (random stragglers).}
Each communication round, $\lfloor 0.5 \times 7 \rfloor + 1 = 4$
clients are randomly designated stragglers; each straggler performs
$E_i \sim \mathrm{Uniform}\{1, \dots, 19\}$ local epochs and the
remaining three perform $E = 20$. Aggregated across the ten seeds,
each client serves as a straggler in approximately $57\%$ of rounds.
This follows the protocol of \citet{li2020fedprox} §5.2 up to two
implementation details: stragglers always perform at least one
optimiser step (the Li et al.\ protocol allows $E_i = 0$), and both
algorithms include partial updates in aggregation (the Li et al.\
canonical FedAvg drops them; we recover the asymmetric protocol
separately in Section~\ref{sec:asymmetric}).

\begin{figure}[!htbp]
\centering
\includegraphics[width=\textwidth]{F6_system_het_schedules.pdf}
\caption{Per-(round, client) local-epoch schedules under the three
system-heterogeneity conditions, reconstructed deterministically from
seed~42. C0 (left) is uniform $E = 20$. C1 (middle) holds clients
C5 and C6 at $E = 5$ in every round. C2 (right) randomly designates
4 of 7 clients per round as stragglers with
$E_i \sim \mathrm{Uniform}\{1, \dots, 19\}$.}
\label{fig:system-het-schedules}
\end{figure}

\subsection{Results on the engineered partition}
\label{sec:system-het-engineered}

Table~\ref{tab:system-het-headline} reports test macro-F1 under the
three conditions on the engineered partition. The mean within-pair
$\Delta$ rises monotonically from $+0.007$ at the no-system-het
baseline (C0) to $+0.013$ under fixed stragglers (C1) to $+0.017$
under random stragglers (C2). FedProx wins 7--8 of 10 paired seeds
in each condition. The pattern is consistent with the canonical
motivation of \citet{li2020fedprox}: as the amount of inexact local
work increases, the proximal regulariser's drift-control mechanism
has more drift to constrain and the within-pair advantage of FedProx
descriptively grows.

% T06 -- System-heterogeneity headline (auto-generated)
\begin{table}[!htbp]
\centering
\caption{System-heterogeneity headline results, all on the Flower runtime with the engineered \texttt{balanced\_paired\_7\_clients} partition and $n=10$ paired seeds. The within-condition $\Delta$ tests H1 (paired Wilcoxon FedProx-vs-FedAvg, two-sided). The between-condition $\Delta-\Delta_{C0}$ tests H2 (paired Wilcoxon on per-seed $\Delta$-of-$\Delta$, Bonferroni-corrected across the two non-baseline conditions).}
\label{tab:system-het-headline}
\small
\begin{tabular}{lcccccc}
\toprule
Condition & FedAvg & FedProx & $\Delta$ & wins & $p_{H1}$ & $p_{H2,\text{Bonf}}$ \\
\midrule
C0 (no system het, baseline) & $0.4964 \pm 0.015$ & $0.5033 \pm 0.021$ & $+0.0069$ & 7/10 & 0.432 & -- \\
C1 (fixed stragglers) & $0.4978 \pm 0.016$ & $0.5111 \pm 0.020$ & $+0.0133$ & 8/10 & 0.160 & 1.000 \\
C2 (random stragglers, primary) & $0.4928 \pm 0.028$ & $0.5102 \pm 0.017$ & $+0.0174$ & 7/10 & 0.232 & 1.000 \\
\bottomrule
\end{tabular}
\end{table}


\paragraph{Per-class breakdown under C2.}
Table~\ref{tab:system-het-per-class} decomposes the C2 result by
class. The pattern matches that observed in the statistical-
heterogeneity section: the FedProx advantage is concentrated on the
melanoma class ($\Delta = +0.085$ macro-F1, FedProx wins 10/10
paired seeds), while the majority class melanocytic nevi shows
essentially zero difference. The proximal-anchor mechanism is
expected to be active only where inter-client drift exists, and the
melanoma class --- held by two specialist clients under the engineered
partition --- is exactly where drift is structurally produced.

% T07 -- Per-class system-het breakdown (auto-generated)
\begin{table}[!htbp]
\centering
\caption{Per-class test macro-F1 under C2 (random stragglers), paired Wilcoxon FedProx-vs-FedAvg with Holm--Bonferroni correction over the seven-class family. Melanoma is the single Holm-corrected significant class ($\Delta = +0.085$, $p_{\text{Holm}} = 0.014$, 10/10 wins).}
\label{tab:system-het-per-class}
\small
\begin{tabular}{lcccccc}
\toprule
Class & FedAvg & FedProx & $\Delta$ & wins & $p_{\text{raw}}$ & $p_{\text{Holm}}$ \\
\midrule
Actinic keratoses & $0.454 \pm 0.060$ & $0.481 \pm 0.038$ & $+0.026$ & 7/10 & 0.432 & 1.000 \\
Basal cell carcinoma & $0.491 \pm 0.057$ & $0.497 \pm 0.023$ & $+0.005$ & 4/10 & 1.000 & 1.000 \\
Benign keratosis & $0.428 \pm 0.066$ & $0.432 \pm 0.045$ & $+0.003$ & 4/10 & 1.000 & 1.000 \\
Dermatofibroma & $0.255 \pm 0.065$ & $0.288 \pm 0.103$ & $+0.034$ & 6/10 & 0.426 & 1.000 \\
Melanoma & $0.233 \pm 0.065$ & $0.318 \pm 0.058$ & $+0.085$ & 10/10 & \textbf{0.002} & \textbf{0.014} \\
Melanocytic nevi & $0.884 \pm 0.005$ & $0.886 \pm 0.004$ & $+0.001$ & 7/10 & 0.557 & 1.000 \\
Vascular lesions & $0.703 \pm 0.066$ & $0.670 \pm 0.039$ & $-0.033$ & 4/10 & 0.131 & 0.785 \\
\bottomrule
\end{tabular}
\end{table}


\paragraph{Straggler-tolerance ratios.}
Table~\ref{tab:system-het-tolerance} reports the ratio
$\rho_a^c = M_a^c / M_a^{C0}$ of each algorithm's mean macro-F1 under
the straggler conditions versus its no-system-het baseline. Both
algorithms maintain $\rho \approx 1$ under both straggler conditions,
indicating that the engineered partition's headline macro-F1 is largely
preserved under the introduction of system heterogeneity. FedProx's
$\rho$ is descriptively higher than FedAvg's in both C1 and C2,
consistent with the per-condition $\Delta$ ordering in
Table~\ref{tab:system-het-headline}.

% T08 -- Straggler-tolerance ratios (auto-generated)
\begin{table}[!htbp]
\centering
\caption{Straggler-tolerance ratios $\rho_a^c = M_a^c / M_a^{C0}$ under each system-heterogeneity condition. Values above one indicate better-than-baseline performance under stragglers, which is implausible on its face and is best read as evidence that the baseline-vs-condition differences are within within-runtime seed noise of approximately $0.005$ macro-F1 at $n = 10$.}
\label{tab:system-het-tolerance}
\small
\begin{tabular}{lcc}
\toprule
Condition & $\rho_{\text{FedAvg}}$ & $\rho_{\text{FedProx}}$ \\
\midrule
C0 (baseline) & $1.000$ & $1.000$ \\
C1 (fixed stragglers) & $1.0028$ & $1.0155$ \\
C2 (random stragglers, primary) & $0.9928$ & $1.0138$ \\
\bottomrule
\end{tabular}
\end{table}


\begin{figure}[!htbp]
\centering
\includegraphics[width=\textwidth]{F7_system_het_convergence.pdf}
\caption{Validation macro-F1 trajectories under the three
system-heterogeneity conditions on each partition design (Flower
runtime, $n = 10$ paired seeds, mean $\pm$ SEM). \textbf{Top row:}
engineered partition; FedProx (dark slate blue) maintains a visible
upward separation from FedAvg (mint green) throughout training in
all three conditions, with the gap widening from C0 to C2.
\textbf{Bottom row:} IID partition; the two algorithms' trajectories
are essentially superimposed in all three conditions, consistent with
the proximal-anchor mechanism's null prediction under i.i.d.\ data.}
\label{fig:system-het-convergence}
\end{figure}

\subsection{Cross-partition control: IID partition}
\label{sec:system-het-iid}

To verify that the system-heterogeneity amplification observed on the
engineered partition reflects the proximal-anchor mechanism specifically
--- rather than a generic regularisation effect --- we ran the same
three conditions on the IID partition, in which every client receives
a uniform random sample of the global distribution and no inter-client
class skew exists. Under the proximal-anchor mechanism, the algorithm
has no drift to constrain on the IID partition and FedProx is
predicted to provide no benefit, irrespective of straggler condition.

Table~\ref{tab:system-het-iid} reports the IID-partition results. The
within-pair $\Delta$ values are small in absolute magnitude under all
three conditions ($-0.006$ at C0, $+0.010$ at C1, $-0.013$ at C2) and
not consistently signed across the three. The two algorithms' macro-F1
distributions overlap substantially on every condition. Whichever
straggler condition is applied, FedProx and FedAvg perform
indistinguishably on the IID partition --- consistent with the
proximal-anchor mechanism's null prediction.

% T12 -- System-het on IID partition (auto-generated)
\begin{table}[!htbp]
\centering
\caption{System-heterogeneity results on the IID partition (\texttt{iid\_7\_clients}) under the Flower runtime with $n = 10$ paired seeds. Values are test macro-F1 mean $\pm$ standard deviation across seeds; the winning algorithm per row is shown in bold. The IID partition serves as a mechanism-null control: with no inter-client class skew, the proximal anchor is theoretically inert.}
\label{tab:system-het-iid}
\small
\begin{tabular}{lcccc}
\toprule
Condition & FedAvg & FedProx & $\Delta$ & Wins \\
\midrule
C0 (no system het) & \textbf{$0.5846 \pm 0.0197$} & $0.5786 \pm 0.0190$ & $-0.0061$ & 5/10 \\
C1 (fixed stragglers) & $0.5741 \pm 0.0188$ & \textbf{$0.5843 \pm 0.0140$} & $+0.0102$ & 6/10 \\
C2 (random stragglers) & \textbf{$0.5879 \pm 0.0112$} & $0.5747 \pm 0.0149$ & $-0.0131$ & 3/10 \\
\bottomrule
\end{tabular}
\end{table}


\paragraph{Cross-partition comparison.}
Table~\ref{tab:system-het-cross-partition} contrasts the
within-pair $\Delta$ values between the two partition designs.
On the engineered partition, $\Delta$ is positive in all three
conditions and grows with straggler severity (mechanism active,
amplification visible). On the IID partition, $\Delta$ is small and
not consistently positive (mechanism inert). The pattern is exactly
the cross-partition signature the proximal-anchor mechanism predicts.

% T13 -- Cross-partition system-het summary (auto-generated)
\begin{table}[!htbp]
\centering
\caption{Cross-partition comparison of the FedProx-vs-FedAvg within-pair difference under each system-heterogeneity condition. The engineered partition contains substantial class skew across clients (every minority class held by two clients); the IID partition contains no inter-client class skew. The proximal anchor's mechanism predicts a positive $\Delta$ on the engineered partition and $\Delta \approx 0$ on IID.}
\label{tab:system-het-cross-partition}
\small
\begin{tabular}{lcccc}
\toprule
& \multicolumn{2}{c}{Engineered partition} & \multicolumn{2}{c}{IID partition} \\
\cmidrule(lr){2-3} \cmidrule(lr){4-5}
Condition & $\Delta$ & Wins & $\Delta$ & Wins \\
\midrule
C0 (no system het) & $+0.0069$ & 7/10 & $-0.0061$ & 5/10 \\
C1 (fixed stragglers) & $+0.0133$ & 8/10 & $+0.0102$ & 6/10 \\
C2 (random stragglers) & $+0.0174$ & 7/10 & $-0.0131$ & 3/10 \\
\bottomrule
\end{tabular}
\end{table}


Figure~\ref{fig:system-het-convergence} shows the trajectories
side-by-side: the top row (engineered) consistently displays a
visible FedProx-FedAvg gap, while the bottom row (IID) shows the
two algorithms' trajectories essentially superimposed. The visual
contrast between the two rows is the most direct illustration of
the mechanism's data-dependence: FedProx helps where there is class
skew, and is inert where there is not.

\subsection{Partial client participation}
\label{sec:partial-participation}

The conditions reported above vary the second canonical FL parameter
of system heterogeneity --- the local-epoch budget $E$ --- while
holding the per-round participation fraction fixed at $C = 1.0$ (every
client participates every round). The third canonical FL parameter,
$C$, is itself an axis of system heterogeneity widely tested in the
literature (\citealt{mcmahan2017fedavg} §3.1 varies
$C \in \{0.1, 0.2, 0.5, 1.0\}$; \citealt{li2020fedprox} §5.3 tests
FedProx under partial participation on Synthetic and FEMNIST;
\citealt{kairouz2021advances} lists $C$ as a defining FL setup
parameter). In real federated deployments, not all clients respond
every round, so the full-participation setting reported above is the
cleanest baseline but the least operationally realistic case. We
additionally report a partial-participation condition at $C = 0.5$
(four of the seven clients sampled uniformly at random per round,
computed as $\max(1, \lfloor 0.5 \times 7 \rfloor) = 4$), holding all
other hyperparameters identical to the engineered C0 baseline.

Table~\ref{tab:partial-participation} reports the result. Under
partial participation the within-pair $\Delta$ is
\TODOhpc{[fill in: $X.XXX$ macro-F1, $Y/10$ paired wins once HPC
results land]}, [direction: descriptively positive / comparable to /
descriptively negative versus] the $C = 1.0$ baseline of $+0.007$
macro-F1. The partial-participation setting introduces additional
aggregation-pool variance per round (fewer clients contributing,
varying subsets across rounds), which on mechanism grounds is
expected to interact with the proximal anchor's drift-control
mechanism in a regime not directly tested by the local-epoch
heterogeneity conditions above. The descriptive result here covers
the second canonical FL system-heterogeneity axis and complements the
local-epoch heterogeneity reported in
Sections~\ref{sec:system-het-engineered}--\ref{sec:system-het-iid}.

\input{T14_partial_participation}

\subsection{Replication of the canonical Li et al.\ §5.2 protocol}
\label{sec:asymmetric}

The H1/H2 tests above use a symmetric inclusion policy in which both
FedAvg and FedProx aggregate over all clients including stragglers.
The canonical asymmetric protocol of \citet{li2020fedprox} §5.2, in
contrast, has FedAvg drop straggler updates entirely while FedProx
includes them as $\gamma$-inexact partial work. The original paper
reports a substantial FedProx advantage under this asymmetric setup.

To verify that this canonical setup reproduces qualitatively on
DermaMNIST, we additionally ran the FedAvg-drop arm on the engineered
partition with random stragglers (the C2 schedule). Combined with the
existing FedProx-include arm from the symmetric C2 data on the same
ten seeds, this gives the within-pair comparison in
Table~\ref{tab:asymmetric-replication}.

% T09 -- Li 2020 §5.2 asymmetric protocol replication (auto-generated)
\begin{table}[!htbp]
\centering
\caption{Replication of the asymmetric straggler protocol of \citet{li2020fedprox} §5.2, in which FedAvg drops straggler updates (its canonical behaviour) while FedProx includes them as $\gamma$-inexact partial work. Values are mean $\pm$ standard deviation of test macro-F1 across $n = 10$ paired seeds on the engineered partition with random stragglers (4 of 7 clients per round). $\Delta_{\mathrm{total}}$ is the within-pair FedProx-include minus FedAvg-drop difference; $p$ is the two-sided paired Wilcoxon $p$-value. \textbf{Bold} marks contrasts significant at $\alpha = 0.05$.}
\label{tab:asymmetric-replication}
\small
\begin{tabular}{lcccc}
\toprule
Algorithm (aggregation policy) & Test macro-F1 & $\Delta_{\mathrm{total}}$ & $p$ (Wilcoxon) & Wins \\
\midrule
FedAvg (drops stragglers) & $0.4423 \pm 0.0192$ & --- & --- & --- \\
FedProx (includes stragglers) & \textbf{$0.5102 \pm 0.0168$} & \textbf{$+0.0679 \pm 0.0175$} & \textbf{$0.002$} & 10/10 \\
\bottomrule
\end{tabular}
\end{table}


The FedProx-include configuration outperforms FedAvg-drop by a mean
within-pair difference of $+0.068$ macro-F1, with FedProx winning
all 10 paired seeds. This reproduces the qualitative finding of
\citet{li2020fedprox} §5.2 on the DermaMNIST task: under the
canonical asymmetric straggler protocol, the FedProx configuration
produces a substantially higher macro-F1 than the FedAvg
configuration.

\subsection{Summary}
\label{sec:system-het-summary}

On the engineered partition the within-pair FedProx-vs-FedAvg
difference rises monotonically with straggler severity (C0:
$+0.007$, C1: $+0.013$, C2: $+0.017$), consistent with the canonical
\citet{li2020fedprox} motivation. The C2 per-class breakdown localises
the advantage to the melanoma class ($\Delta = +0.085$ macro-F1, 10/10
paired wins). Under partial client participation ($C = 0.5$,
Section~\ref{sec:partial-participation}) the descriptive result is
\TODOhpc{[fill in once HPC data lands]}. The canonical asymmetric
protocol of \citet{li2020fedprox} §5.2 reproduces qualitatively:
$\Delta_{\text{total}} = +0.068$ macro-F1 with FedProx winning all 10
paired seeds. On the IID partition the within-pair $\Delta$ is small
and not consistently positive under any of the three straggler
conditions, in line with the proximal-anchor mechanism's prediction
that the regulariser is inert when no inter-client class skew exists.

\paragraph{Scope and limitations.}
The within-pair $\Delta$ magnitudes on the engineered partition
($+0.007$--$+0.017$ macro-F1) are smaller than the effect sizes
originally reported in \citet{li2020fedprox} on Synthetic, FEMNIST,
Shakespeare, and Sentiment140 benchmarks. This is consistent with
the small absolute magnitude observed in the statistical-heterogeneity
arm under the same partition and runtime, and is plausibly attributable
to the smaller dataset, more challenging architecture, and
medical-imaging difficulty discussed in
Section~\ref{sec:headline-summary}. The C1 fixed-straggler condition
on the engineered partition is structurally coupled with the clinical
role assignment (C5 holds the melanoma/vascular pair); its descriptive
result should be read with that confound in mind. The asymmetric
protocol of Section~\ref{sec:asymmetric} provides a single point of
comparison against the canonical Li et al.\ setup but is not a primary
analysis; it is included to demonstrate qualitative reproducibility of
their §5.2 result on DermaMNIST. The straggler schedules are
deterministic functions of the seed and have been verified by direct
reconstruction from the JSON provenance --- the empirical per-client
straggler frequency, $E_i$ range, and round counts match the schedule
specifications exactly.

% ==============================================================================

\section{System Heterogeneity}
\label{sec:results-system-het}

System heterogeneity refers to differences in the local computational
budget that participating clients contribute per communication round.
In production federated deployments this is the dominant operational
challenge: institutions run on heterogeneous hardware, on-call training
windows differ by site, and not all clients can complete the same
amount of local work within a round's wall-clock deadline. The proximal
term in FedProx was originally motivated for this regime by
\citet{li2020fedprox}: clients may submit $\gamma$-inexact updates ---
partial work that does not fully minimise the local objective --- while
the proximal anchor $(\mu/2)\lVert w - w^t \rVert^2$ provably preserves
convergence by preventing inexact updates from drifting too far from
the broadcast iterate $w^t$ (Theorem~4).

This section reports the FedProx-vs-FedAvg comparison under three
straggler conditions on the engineered \texttt{balanced\_paired\_7\_clients}
partition (the same partition used as the primary in
Section~\ref{sec:results-statistical-het}) and additionally on the IID
partition (\texttt{iid\_7\_clients}) as a mechanism control. We follow
the same descriptive reporting convention as the
statistical-heterogeneity section: mean $\pm$ SD across paired seeds,
visual curve comparisons, bolded best results.

\paragraph{Terminology.}
We use ``system heterogeneity'' as shorthand for
\emph{computational-work heterogeneity}: variability in the number of
local SGD steps each client performs per round. The manipulation acts
only on the algorithm via heterogeneous local-epoch budgets $\tau_i$;
we do not simulate hardware speed, network latency, or device
availability.

\subsection{Experimental design and condition motivation}
\label{sec:system-het-conditions}

Three system-heterogeneity conditions are evaluated. Each schedule is
deterministic given the seed (see \texttt{fl/system\_het.py}), so
paired FedAvg/FedProx runs share an identical per-(round, client)
local-epoch assignment and the within-pair $\Delta$ reflects only the
algorithmic difference. The schedules under seed~42 are visualised in
Figure~\ref{fig:system-het-schedules}.

\paragraph{C0 (no system heterogeneity, baseline).}
Every client performs $E = 20$ local epochs every round. The design
is identical to the headline statistical-heterogeneity setup, but
re-executed in the Flower runtime so that all three system-het
conditions share a single runtime. C0 provides the within-runtime
baseline against which the system-het amplification can be assessed.

\paragraph{C1 (fixed stragglers).}
Two designated clients (C5 and C6) always perform
$E_{\text{straggler}} = 5$ local epochs; the remaining five clients
perform the full $E = 20$. Under the engineered partition, C5 holds
the melanoma/vascular specialty pair and C6 is the nevi-only
generalist, so the C1 design is structurally coupled with the
clinical-role assignment under that partition. Under the IID partition
the choice of (C5, C6) is arbitrary (clients are exchangeable), so the
IID C1 isolates the pure compute-heterogeneity effect more cleanly.

\paragraph{C2 (random stragglers).}
Each communication round, $\lfloor 0.5 \times 7 \rfloor + 1 = 4$
clients are randomly designated stragglers; each straggler performs
$E_i \sim \mathrm{Uniform}\{1, \dots, 19\}$ local epochs and the
remaining three perform $E = 20$. Aggregated across the ten seeds,
each client serves as a straggler in approximately $57\%$ of rounds.
This follows the protocol of \citet{li2020fedprox} §5.2 up to two
implementation details: stragglers always perform at least one
optimiser step (the Li et al.\ protocol allows $E_i = 0$), and both
algorithms include partial updates in aggregation (the Li et al.\
canonical FedAvg drops them; we recover the asymmetric protocol
separately in Section~\ref{sec:asymmetric}).

\begin{figure}[!htbp]
\centering
\includegraphics[width=\textwidth]{F6_system_het_schedules.pdf}
\caption{Per-(round, client) local-epoch schedules under the three
system-heterogeneity conditions, reconstructed deterministically from
seed~42. C0 (left) is uniform $E = 20$. C1 (middle) holds clients
C5 and C6 at $E = 5$ in every round. C2 (right) randomly designates
4 of 7 clients per round as stragglers with
$E_i \sim \mathrm{Uniform}\{1, \dots, 19\}$.}
\label{fig:system-het-schedules}
\end{figure}

\subsection{Results on the engineered partition}
\label{sec:system-het-engineered}

Table~\ref{tab:system-het-headline} reports test macro-F1 under the
three conditions on the engineered partition. The mean within-pair
$\Delta$ rises monotonically from $+0.007$ at the no-system-het
baseline (C0) to $+0.013$ under fixed stragglers (C1) to $+0.017$
under random stragglers (C2). FedProx wins 7--8 of 10 paired seeds
in each condition. The pattern is consistent with the canonical
motivation of \citet{li2020fedprox}: as the amount of inexact local
work increases, the proximal regulariser's drift-control mechanism
has more drift to constrain and the within-pair advantage of FedProx
descriptively grows.

% T06 -- System-heterogeneity headline (auto-generated)
\begin{table}[!htbp]
\centering
\caption{System-heterogeneity headline results, all on the Flower runtime with the engineered \texttt{balanced\_paired\_7\_clients} partition and $n=10$ paired seeds. The within-condition $\Delta$ tests H1 (paired Wilcoxon FedProx-vs-FedAvg, two-sided). The between-condition $\Delta-\Delta_{C0}$ tests H2 (paired Wilcoxon on per-seed $\Delta$-of-$\Delta$, Bonferroni-corrected across the two non-baseline conditions).}
\label{tab:system-het-headline}
\small
\begin{tabular}{lcccccc}
\toprule
Condition & FedAvg & FedProx & $\Delta$ & wins & $p_{H1}$ & $p_{H2,\text{Bonf}}$ \\
\midrule
C0 (no system het, baseline) & $0.4964 \pm 0.015$ & $0.5033 \pm 0.021$ & $+0.0069$ & 7/10 & 0.432 & -- \\
C1 (fixed stragglers) & $0.4978 \pm 0.016$ & $0.5111 \pm 0.020$ & $+0.0133$ & 8/10 & 0.160 & 1.000 \\
C2 (random stragglers, primary) & $0.4928 \pm 0.028$ & $0.5102 \pm 0.017$ & $+0.0174$ & 7/10 & 0.232 & 1.000 \\
\bottomrule
\end{tabular}
\end{table}


\paragraph{Per-class breakdown under C2.}
Table~\ref{tab:system-het-per-class} decomposes the C2 result by
class. The pattern matches that observed in the statistical-
heterogeneity section: the FedProx advantage is concentrated on the
melanoma class ($\Delta = +0.085$ macro-F1, FedProx wins 10/10
paired seeds), while the majority class melanocytic nevi shows
essentially zero difference. The proximal-anchor mechanism is
expected to be active only where inter-client drift exists, and the
melanoma class --- held by two specialist clients under the engineered
partition --- is exactly where drift is structurally produced.

% T07 -- Per-class system-het breakdown (auto-generated)
\begin{table}[!htbp]
\centering
\caption{Per-class test macro-F1 under C2 (random stragglers), paired Wilcoxon FedProx-vs-FedAvg with Holm--Bonferroni correction over the seven-class family. Melanoma is the single Holm-corrected significant class ($\Delta = +0.085$, $p_{\text{Holm}} = 0.014$, 10/10 wins).}
\label{tab:system-het-per-class}
\small
\begin{tabular}{lcccccc}
\toprule
Class & FedAvg & FedProx & $\Delta$ & wins & $p_{\text{raw}}$ & $p_{\text{Holm}}$ \\
\midrule
Actinic keratoses & $0.454 \pm 0.060$ & $0.481 \pm 0.038$ & $+0.026$ & 7/10 & 0.432 & 1.000 \\
Basal cell carcinoma & $0.491 \pm 0.057$ & $0.497 \pm 0.023$ & $+0.005$ & 4/10 & 1.000 & 1.000 \\
Benign keratosis & $0.428 \pm 0.066$ & $0.432 \pm 0.045$ & $+0.003$ & 4/10 & 1.000 & 1.000 \\
Dermatofibroma & $0.255 \pm 0.065$ & $0.288 \pm 0.103$ & $+0.034$ & 6/10 & 0.426 & 1.000 \\
Melanoma & $0.233 \pm 0.065$ & $0.318 \pm 0.058$ & $+0.085$ & 10/10 & \textbf{0.002} & \textbf{0.014} \\
Melanocytic nevi & $0.884 \pm 0.005$ & $0.886 \pm 0.004$ & $+0.001$ & 7/10 & 0.557 & 1.000 \\
Vascular lesions & $0.703 \pm 0.066$ & $0.670 \pm 0.039$ & $-0.033$ & 4/10 & 0.131 & 0.785 \\
\bottomrule
\end{tabular}
\end{table}


\paragraph{Straggler-tolerance ratios.}
Table~\ref{tab:system-het-tolerance} reports the ratio
$\rho_a^c = M_a^c / M_a^{C0}$ of each algorithm's mean macro-F1 under
the straggler conditions versus its no-system-het baseline. Both
algorithms maintain $\rho \approx 1$ under both straggler conditions,
indicating that the engineered partition's headline macro-F1 is largely
preserved under the introduction of system heterogeneity. FedProx's
$\rho$ is descriptively higher than FedAvg's in both C1 and C2,
consistent with the per-condition $\Delta$ ordering in
Table~\ref{tab:system-het-headline}.

% T08 -- Straggler-tolerance ratios (auto-generated)
\begin{table}[!htbp]
\centering
\caption{Straggler-tolerance ratios $\rho_a^c = M_a^c / M_a^{C0}$ under each system-heterogeneity condition. Values above one indicate better-than-baseline performance under stragglers, which is implausible on its face and is best read as evidence that the baseline-vs-condition differences are within within-runtime seed noise of approximately $0.005$ macro-F1 at $n = 10$.}
\label{tab:system-het-tolerance}
\small
\begin{tabular}{lcc}
\toprule
Condition & $\rho_{\text{FedAvg}}$ & $\rho_{\text{FedProx}}$ \\
\midrule
C0 (baseline) & $1.000$ & $1.000$ \\
C1 (fixed stragglers) & $1.0028$ & $1.0155$ \\
C2 (random stragglers, primary) & $0.9928$ & $1.0138$ \\
\bottomrule
\end{tabular}
\end{table}


\begin{figure}[!htbp]
\centering
\includegraphics[width=\textwidth]{F7_system_het_convergence.pdf}
\caption{Validation macro-F1 trajectories under the three
system-heterogeneity conditions on each partition design (Flower
runtime, $n = 10$ paired seeds, mean $\pm$ SEM). \textbf{Top row:}
engineered partition; FedProx (dark slate blue) maintains a visible
upward separation from FedAvg (mint green) throughout training in
all three conditions, with the gap widening from C0 to C2.
\textbf{Bottom row:} IID partition; the two algorithms' trajectories
are essentially superimposed in all three conditions, consistent with
the proximal-anchor mechanism's null prediction under i.i.d.\ data.}
\label{fig:system-het-convergence}
\end{figure}

\subsection{Cross-partition control: IID partition}
\label{sec:system-het-iid}

To verify that the system-heterogeneity amplification observed on the
engineered partition reflects the proximal-anchor mechanism specifically
--- rather than a generic regularisation effect --- we ran the same
three conditions on the IID partition, in which every client receives
a uniform random sample of the global distribution and no inter-client
class skew exists. Under the proximal-anchor mechanism, the algorithm
has no drift to constrain on the IID partition and FedProx is
predicted to provide no benefit, irrespective of straggler condition.

Table~\ref{tab:system-het-iid} reports the IID-partition results. The
within-pair $\Delta$ values are small in absolute magnitude under all
three conditions ($-0.006$ at C0, $+0.010$ at C1, $-0.013$ at C2) and
not consistently signed across the three. The two algorithms' macro-F1
distributions overlap substantially on every condition. Whichever
straggler condition is applied, FedProx and FedAvg perform
indistinguishably on the IID partition --- consistent with the
proximal-anchor mechanism's null prediction.

% T12 -- System-het on IID partition (auto-generated)
\begin{table}[!htbp]
\centering
\caption{System-heterogeneity results on the IID partition (\texttt{iid\_7\_clients}) under the Flower runtime with $n = 10$ paired seeds. Values are test macro-F1 mean $\pm$ standard deviation across seeds; the winning algorithm per row is shown in bold. The IID partition serves as a mechanism-null control: with no inter-client class skew, the proximal anchor is theoretically inert.}
\label{tab:system-het-iid}
\small
\begin{tabular}{lcccc}
\toprule
Condition & FedAvg & FedProx & $\Delta$ & Wins \\
\midrule
C0 (no system het) & \textbf{$0.5846 \pm 0.0197$} & $0.5786 \pm 0.0190$ & $-0.0061$ & 5/10 \\
C1 (fixed stragglers) & $0.5741 \pm 0.0188$ & \textbf{$0.5843 \pm 0.0140$} & $+0.0102$ & 6/10 \\
C2 (random stragglers) & \textbf{$0.5879 \pm 0.0112$} & $0.5747 \pm 0.0149$ & $-0.0131$ & 3/10 \\
\bottomrule
\end{tabular}
\end{table}


\paragraph{Cross-partition comparison.}
Table~\ref{tab:system-het-cross-partition} contrasts the
within-pair $\Delta$ values between the two partition designs.
On the engineered partition, $\Delta$ is positive in all three
conditions and grows with straggler severity (mechanism active,
amplification visible). On the IID partition, $\Delta$ is small and
not consistently positive (mechanism inert). The pattern is exactly
the cross-partition signature the proximal-anchor mechanism predicts.

% T13 -- Cross-partition system-het summary (auto-generated)
\begin{table}[!htbp]
\centering
\caption{Cross-partition comparison of the FedProx-vs-FedAvg within-pair difference under each system-heterogeneity condition. The engineered partition contains substantial class skew across clients (every minority class held by two clients); the IID partition contains no inter-client class skew. The proximal anchor's mechanism predicts a positive $\Delta$ on the engineered partition and $\Delta \approx 0$ on IID.}
\label{tab:system-het-cross-partition}
\small
\begin{tabular}{lcccc}
\toprule
& \multicolumn{2}{c}{Engineered partition} & \multicolumn{2}{c}{IID partition} \\
\cmidrule(lr){2-3} \cmidrule(lr){4-5}
Condition & $\Delta$ & Wins & $\Delta$ & Wins \\
\midrule
C0 (no system het) & $+0.0069$ & 7/10 & $-0.0061$ & 5/10 \\
C1 (fixed stragglers) & $+0.0133$ & 8/10 & $+0.0102$ & 6/10 \\
C2 (random stragglers) & $+0.0174$ & 7/10 & $-0.0131$ & 3/10 \\
\bottomrule
\end{tabular}
\end{table}


Figure~\ref{fig:system-het-convergence} shows the trajectories
side-by-side: the top row (engineered) consistently displays a
visible FedProx-FedAvg gap, while the bottom row (IID) shows the
two algorithms' trajectories essentially superimposed. The visual
contrast between the two rows is the most direct illustration of
the mechanism's data-dependence: FedProx helps where there is class
skew, and is inert where there is not.

\subsection{Partial client participation}
\label{sec:partial-participation}

The conditions reported above vary the second canonical FL parameter
of system heterogeneity --- the local-epoch budget $E$ --- while
holding the per-round participation fraction fixed at $C = 1.0$ (every
client participates every round). The third canonical FL parameter,
$C$, is itself an axis of system heterogeneity widely tested in the
literature (\citealt{mcmahan2017fedavg} §3.1 varies
$C \in \{0.1, 0.2, 0.5, 1.0\}$; \citealt{li2020fedprox} §5.3 tests
FedProx under partial participation on Synthetic and FEMNIST;
\citealt{kairouz2021advances} lists $C$ as a defining FL setup
parameter). In real federated deployments, not all clients respond
every round, so the full-participation setting reported above is the
cleanest baseline but the least operationally realistic case. We
additionally report a partial-participation condition at $C = 0.5$
(four of the seven clients sampled uniformly at random per round,
computed as $\max(1, \lfloor 0.5 \times 7 \rfloor) = 4$), holding all
other hyperparameters identical to the engineered C0 baseline.

Table~\ref{tab:partial-participation} reports the result. Under
partial participation the within-pair $\Delta$ is
\TODOhpc{[fill in: $X.XXX$ macro-F1, $Y/10$ paired wins once HPC
results land]}, [direction: descriptively positive / comparable to /
descriptively negative versus] the $C = 1.0$ baseline of $+0.007$
macro-F1. The partial-participation setting introduces additional
aggregation-pool variance per round (fewer clients contributing,
varying subsets across rounds), which on mechanism grounds is
expected to interact with the proximal anchor's drift-control
mechanism in a regime not directly tested by the local-epoch
heterogeneity conditions above. The descriptive result here covers
the second canonical FL system-heterogeneity axis and complements the
local-epoch heterogeneity reported in
Sections~\ref{sec:system-het-engineered}--\ref{sec:system-het-iid}.

\input{T14_partial_participation}

\subsection{Replication of the canonical Li et al.\ §5.2 protocol}
\label{sec:asymmetric}

The H1/H2 tests above use a symmetric inclusion policy in which both
FedAvg and FedProx aggregate over all clients including stragglers.
The canonical asymmetric protocol of \citet{li2020fedprox} §5.2, in
contrast, has FedAvg drop straggler updates entirely while FedProx
includes them as $\gamma$-inexact partial work. The original paper
reports a substantial FedProx advantage under this asymmetric setup.

To verify that this canonical setup reproduces qualitatively on
DermaMNIST, we additionally ran the FedAvg-drop arm on the engineered
partition with random stragglers (the C2 schedule). Combined with the
existing FedProx-include arm from the symmetric C2 data on the same
ten seeds, this gives the within-pair comparison in
Table~\ref{tab:asymmetric-replication}.

% T09 -- Li 2020 §5.2 asymmetric protocol replication (auto-generated)
\begin{table}[!htbp]
\centering
\caption{Replication of the asymmetric straggler protocol of \citet{li2020fedprox} §5.2, in which FedAvg drops straggler updates (its canonical behaviour) while FedProx includes them as $\gamma$-inexact partial work. Values are mean $\pm$ standard deviation of test macro-F1 across $n = 10$ paired seeds on the engineered partition with random stragglers (4 of 7 clients per round). $\Delta_{\mathrm{total}}$ is the within-pair FedProx-include minus FedAvg-drop difference; $p$ is the two-sided paired Wilcoxon $p$-value. \textbf{Bold} marks contrasts significant at $\alpha = 0.05$.}
\label{tab:asymmetric-replication}
\small
\begin{tabular}{lcccc}
\toprule
Algorithm (aggregation policy) & Test macro-F1 & $\Delta_{\mathrm{total}}$ & $p$ (Wilcoxon) & Wins \\
\midrule
FedAvg (drops stragglers) & $0.4423 \pm 0.0192$ & --- & --- & --- \\
FedProx (includes stragglers) & \textbf{$0.5102 \pm 0.0168$} & \textbf{$+0.0679 \pm 0.0175$} & \textbf{$0.002$} & 10/10 \\
\bottomrule
\end{tabular}
\end{table}


The FedProx-include configuration outperforms FedAvg-drop by a mean
within-pair difference of $+0.068$ macro-F1, with FedProx winning
all 10 paired seeds. This reproduces the qualitative finding of
\citet{li2020fedprox} §5.2 on the DermaMNIST task: under the
canonical asymmetric straggler protocol, the FedProx configuration
produces a substantially higher macro-F1 than the FedAvg
configuration.

\subsection{Summary}
\label{sec:system-het-summary}

On the engineered partition the within-pair FedProx-vs-FedAvg
difference rises monotonically with straggler severity (C0:
$+0.007$, C1: $+0.013$, C2: $+0.017$), consistent with the canonical
\citet{li2020fedprox} motivation. The C2 per-class breakdown localises
the advantage to the melanoma class ($\Delta = +0.085$ macro-F1, 10/10
paired wins). Under partial client participation ($C = 0.5$,
Section~\ref{sec:partial-participation}) the descriptive result is
\TODOhpc{[fill in once HPC data lands]}. The canonical asymmetric
protocol of \citet{li2020fedprox} §5.2 reproduces qualitatively:
$\Delta_{\text{total}} = +0.068$ macro-F1 with FedProx winning all 10
paired seeds. On the IID partition the within-pair $\Delta$ is small
and not consistently positive under any of the three straggler
conditions, in line with the proximal-anchor mechanism's prediction
that the regulariser is inert when no inter-client class skew exists.

\paragraph{Scope and limitations.}
The within-pair $\Delta$ magnitudes on the engineered partition
($+0.007$--$+0.017$ macro-F1) are smaller than the effect sizes
originally reported in \citet{li2020fedprox} on Synthetic, FEMNIST,
Shakespeare, and Sentiment140 benchmarks. This is consistent with
the small absolute magnitude observed in the statistical-heterogeneity
arm under the same partition and runtime, and is plausibly attributable
to the smaller dataset, more challenging architecture, and
medical-imaging difficulty discussed in
Section~\ref{sec:headline-summary}. The C1 fixed-straggler condition
on the engineered partition is structurally coupled with the clinical
role assignment (C5 holds the melanoma/vascular pair); its descriptive
result should be read with that confound in mind. The asymmetric
protocol of Section~\ref{sec:asymmetric} provides a single point of
comparison against the canonical Li et al.\ setup but is not a primary
analysis; it is included to demonstrate qualitative reproducibility of
their §5.2 result on DermaMNIST. The straggler schedules are
deterministic functions of the seed and have been verified by direct
reconstruction from the JSON provenance --- the empirical per-client
straggler frequency, $E_i$ range, and round counts match the schedule
specifications exactly.

% ==============================================================================

\section{System Heterogeneity}
\label{sec:results-system-het}

System heterogeneity refers to differences in the local computational
budget that participating clients contribute per communication round.
In production federated deployments this is the dominant operational
challenge: institutions run on heterogeneous hardware, on-call training
windows differ by site, and not all clients can complete the same
amount of local work within a round's wall-clock deadline. The proximal
term in FedProx was originally motivated for this regime by
\citet{li2020fedprox}: clients may submit $\gamma$-inexact updates ---
partial work that does not fully minimise the local objective --- while
the proximal anchor $(\mu/2)\lVert w - w^t \rVert^2$ provably preserves
convergence by preventing inexact updates from drifting too far from
the broadcast iterate $w^t$ (Theorem~4).

This section reports the FedProx-vs-FedAvg comparison under three
straggler conditions on the engineered \texttt{balanced\_paired\_7\_clients}
partition (the same partition used as the primary in
Section~\ref{sec:results-statistical-het}) and additionally on the IID
partition (\texttt{iid\_7\_clients}) as a mechanism control. We follow
the same descriptive reporting convention as the
statistical-heterogeneity section: mean $\pm$ SD across paired seeds,
visual curve comparisons, bolded best results.

\paragraph{Terminology.}
We use ``system heterogeneity'' as shorthand for
\emph{computational-work heterogeneity}: variability in the number of
local SGD steps each client performs per round. The manipulation acts
only on the algorithm via heterogeneous local-epoch budgets $\tau_i$;
we do not simulate hardware speed, network latency, or device
availability.

\subsection{Experimental design and condition motivation}
\label{sec:system-het-conditions}

Three system-heterogeneity conditions are evaluated. Each schedule is
deterministic given the seed (see \texttt{fl/system\_het.py}), so
paired FedAvg/FedProx runs share an identical per-(round, client)
local-epoch assignment and the within-pair $\Delta$ reflects only the
algorithmic difference. The schedules under seed~42 are visualised in
Figure~\ref{fig:system-het-schedules}.

\paragraph{C0 (no system heterogeneity, baseline).}
Every client performs $E = 20$ local epochs every round. The design
is identical to the headline statistical-heterogeneity setup, but
re-executed in the Flower runtime so that all three system-het
conditions share a single runtime. C0 provides the within-runtime
baseline against which the system-het amplification can be assessed.

\paragraph{C1 (fixed stragglers).}
Two designated clients (C5 and C6) always perform
$E_{\text{straggler}} = 5$ local epochs; the remaining five clients
perform the full $E = 20$. Under the engineered partition, C5 holds
the melanoma/vascular specialty pair and C6 is the nevi-only
generalist, so the C1 design is structurally coupled with the
clinical-role assignment under that partition. Under the IID partition
the choice of (C5, C6) is arbitrary (clients are exchangeable), so the
IID C1 isolates the pure compute-heterogeneity effect more cleanly.

\paragraph{C2 (random stragglers).}
Each communication round, $\lfloor 0.5 \times 7 \rfloor + 1 = 4$
clients are randomly designated stragglers; each straggler performs
$E_i \sim \mathrm{Uniform}\{1, \dots, 19\}$ local epochs and the
remaining three perform $E = 20$. Aggregated across the ten seeds,
each client serves as a straggler in approximately $57\%$ of rounds.
This follows the protocol of \citet{li2020fedprox} §5.2 up to two
implementation details: stragglers always perform at least one
optimiser step (the Li et al.\ protocol allows $E_i = 0$), and both
algorithms include partial updates in aggregation (the Li et al.\
canonical FedAvg drops them; we recover the asymmetric protocol
separately in Section~\ref{sec:asymmetric}).

\begin{figure}[!htbp]
\centering
\includegraphics[width=\textwidth]{F6_system_het_schedules.pdf}
\caption{Per-(round, client) local-epoch schedules under the three
system-heterogeneity conditions, reconstructed deterministically from
seed~42. C0 (left) is uniform $E = 20$. C1 (middle) holds clients
C5 and C6 at $E = 5$ in every round. C2 (right) randomly designates
4 of 7 clients per round as stragglers with
$E_i \sim \mathrm{Uniform}\{1, \dots, 19\}$.}
\label{fig:system-het-schedules}
\end{figure}

\subsection{Results on the engineered partition}
\label{sec:system-het-engineered}

Table~\ref{tab:system-het-headline} reports test macro-F1 under the
three conditions on the engineered partition. The mean within-pair
$\Delta$ rises monotonically from $+0.007$ at the no-system-het
baseline (C0) to $+0.013$ under fixed stragglers (C1) to $+0.017$
under random stragglers (C2). FedProx wins 7--8 of 10 paired seeds
in each condition. The pattern is consistent with the canonical
motivation of \citet{li2020fedprox}: as the amount of inexact local
work increases, the proximal regulariser's drift-control mechanism
has more drift to constrain and the within-pair advantage of FedProx
descriptively grows.

% T06 -- System-heterogeneity headline (auto-generated)
\begin{table}[!htbp]
\centering
\caption{System-heterogeneity headline results, all on the Flower runtime with the engineered \texttt{balanced\_paired\_7\_clients} partition and $n=10$ paired seeds. The within-condition $\Delta$ tests H1 (paired Wilcoxon FedProx-vs-FedAvg, two-sided). The between-condition $\Delta-\Delta_{C0}$ tests H2 (paired Wilcoxon on per-seed $\Delta$-of-$\Delta$, Bonferroni-corrected across the two non-baseline conditions).}
\label{tab:system-het-headline}
\small
\begin{tabular}{lcccccc}
\toprule
Condition & FedAvg & FedProx & $\Delta$ & wins & $p_{H1}$ & $p_{H2,\text{Bonf}}$ \\
\midrule
C0 (no system het, baseline) & $0.4964 \pm 0.015$ & $0.5033 \pm 0.021$ & $+0.0069$ & 7/10 & 0.432 & -- \\
C1 (fixed stragglers) & $0.4978 \pm 0.016$ & $0.5111 \pm 0.020$ & $+0.0133$ & 8/10 & 0.160 & 1.000 \\
C2 (random stragglers, primary) & $0.4928 \pm 0.028$ & $0.5102 \pm 0.017$ & $+0.0174$ & 7/10 & 0.232 & 1.000 \\
\bottomrule
\end{tabular}
\end{table}


\paragraph{Per-class breakdown under C2.}
Table~\ref{tab:system-het-per-class} decomposes the C2 result by
class. The pattern matches that observed in the statistical-
heterogeneity section: the FedProx advantage is concentrated on the
melanoma class ($\Delta = +0.085$ macro-F1, FedProx wins 10/10
paired seeds), while the majority class melanocytic nevi shows
essentially zero difference. The proximal-anchor mechanism is
expected to be active only where inter-client drift exists, and the
melanoma class --- held by two specialist clients under the engineered
partition --- is exactly where drift is structurally produced.

% T07 -- Per-class system-het breakdown (auto-generated)
\begin{table}[!htbp]
\centering
\caption{Per-class test macro-F1 under C2 (random stragglers), paired Wilcoxon FedProx-vs-FedAvg with Holm--Bonferroni correction over the seven-class family. Melanoma is the single Holm-corrected significant class ($\Delta = +0.085$, $p_{\text{Holm}} = 0.014$, 10/10 wins).}
\label{tab:system-het-per-class}
\small
\begin{tabular}{lcccccc}
\toprule
Class & FedAvg & FedProx & $\Delta$ & wins & $p_{\text{raw}}$ & $p_{\text{Holm}}$ \\
\midrule
Actinic keratoses & $0.454 \pm 0.060$ & $0.481 \pm 0.038$ & $+0.026$ & 7/10 & 0.432 & 1.000 \\
Basal cell carcinoma & $0.491 \pm 0.057$ & $0.497 \pm 0.023$ & $+0.005$ & 4/10 & 1.000 & 1.000 \\
Benign keratosis & $0.428 \pm 0.066$ & $0.432 \pm 0.045$ & $+0.003$ & 4/10 & 1.000 & 1.000 \\
Dermatofibroma & $0.255 \pm 0.065$ & $0.288 \pm 0.103$ & $+0.034$ & 6/10 & 0.426 & 1.000 \\
Melanoma & $0.233 \pm 0.065$ & $0.318 \pm 0.058$ & $+0.085$ & 10/10 & \textbf{0.002} & \textbf{0.014} \\
Melanocytic nevi & $0.884 \pm 0.005$ & $0.886 \pm 0.004$ & $+0.001$ & 7/10 & 0.557 & 1.000 \\
Vascular lesions & $0.703 \pm 0.066$ & $0.670 \pm 0.039$ & $-0.033$ & 4/10 & 0.131 & 0.785 \\
\bottomrule
\end{tabular}
\end{table}


\paragraph{Straggler-tolerance ratios.}
Table~\ref{tab:system-het-tolerance} reports the ratio
$\rho_a^c = M_a^c / M_a^{C0}$ of each algorithm's mean macro-F1 under
the straggler conditions versus its no-system-het baseline. Both
algorithms maintain $\rho \approx 1$ under both straggler conditions,
indicating that the engineered partition's headline macro-F1 is largely
preserved under the introduction of system heterogeneity. FedProx's
$\rho$ is descriptively higher than FedAvg's in both C1 and C2,
consistent with the per-condition $\Delta$ ordering in
Table~\ref{tab:system-het-headline}.

% T08 -- Straggler-tolerance ratios (auto-generated)
\begin{table}[!htbp]
\centering
\caption{Straggler-tolerance ratios $\rho_a^c = M_a^c / M_a^{C0}$ under each system-heterogeneity condition. Values above one indicate better-than-baseline performance under stragglers, which is implausible on its face and is best read as evidence that the baseline-vs-condition differences are within within-runtime seed noise of approximately $0.005$ macro-F1 at $n = 10$.}
\label{tab:system-het-tolerance}
\small
\begin{tabular}{lcc}
\toprule
Condition & $\rho_{\text{FedAvg}}$ & $\rho_{\text{FedProx}}$ \\
\midrule
C0 (baseline) & $1.000$ & $1.000$ \\
C1 (fixed stragglers) & $1.0028$ & $1.0155$ \\
C2 (random stragglers, primary) & $0.9928$ & $1.0138$ \\
\bottomrule
\end{tabular}
\end{table}


\begin{figure}[!htbp]
\centering
\includegraphics[width=\textwidth]{F7_system_het_convergence.pdf}
\caption{Validation macro-F1 trajectories under the three
system-heterogeneity conditions on each partition design (Flower
runtime, $n = 10$ paired seeds, mean $\pm$ SEM). \textbf{Top row:}
engineered partition; FedProx (dark slate blue) maintains a visible
upward separation from FedAvg (mint green) throughout training in
all three conditions, with the gap widening from C0 to C2.
\textbf{Bottom row:} IID partition; the two algorithms' trajectories
are essentially superimposed in all three conditions, consistent with
the proximal-anchor mechanism's null prediction under i.i.d.\ data.}
\label{fig:system-het-convergence}
\end{figure}

\subsection{Cross-partition control: IID partition}
\label{sec:system-het-iid}

To verify that the system-heterogeneity amplification observed on the
engineered partition reflects the proximal-anchor mechanism specifically
--- rather than a generic regularisation effect --- we ran the same
three conditions on the IID partition, in which every client receives
a uniform random sample of the global distribution and no inter-client
class skew exists. Under the proximal-anchor mechanism, the algorithm
has no drift to constrain on the IID partition and FedProx is
predicted to provide no benefit, irrespective of straggler condition.

Table~\ref{tab:system-het-iid} reports the IID-partition results. The
within-pair $\Delta$ values are small in absolute magnitude under all
three conditions ($-0.006$ at C0, $+0.010$ at C1, $-0.013$ at C2) and
not consistently signed across the three. The two algorithms' macro-F1
distributions overlap substantially on every condition. Whichever
straggler condition is applied, FedProx and FedAvg perform
indistinguishably on the IID partition --- consistent with the
proximal-anchor mechanism's null prediction.

% T12 -- System-het on IID partition (auto-generated)
\begin{table}[!htbp]
\centering
\caption{System-heterogeneity results on the IID partition (\texttt{iid\_7\_clients}) under the Flower runtime with $n = 10$ paired seeds. Values are test macro-F1 mean $\pm$ standard deviation across seeds; the winning algorithm per row is shown in bold. The IID partition serves as a mechanism-null control: with no inter-client class skew, the proximal anchor is theoretically inert.}
\label{tab:system-het-iid}
\small
\begin{tabular}{lcccc}
\toprule
Condition & FedAvg & FedProx & $\Delta$ & Wins \\
\midrule
C0 (no system het) & \textbf{$0.5846 \pm 0.0197$} & $0.5786 \pm 0.0190$ & $-0.0061$ & 5/10 \\
C1 (fixed stragglers) & $0.5741 \pm 0.0188$ & \textbf{$0.5843 \pm 0.0140$} & $+0.0102$ & 6/10 \\
C2 (random stragglers) & \textbf{$0.5879 \pm 0.0112$} & $0.5747 \pm 0.0149$ & $-0.0131$ & 3/10 \\
\bottomrule
\end{tabular}
\end{table}


\paragraph{Cross-partition comparison.}
Table~\ref{tab:system-het-cross-partition} contrasts the
within-pair $\Delta$ values between the two partition designs.
On the engineered partition, $\Delta$ is positive in all three
conditions and grows with straggler severity (mechanism active,
amplification visible). On the IID partition, $\Delta$ is small and
not consistently positive (mechanism inert). The pattern is exactly
the cross-partition signature the proximal-anchor mechanism predicts.

% T13 -- Cross-partition system-het summary (auto-generated)
\begin{table}[!htbp]
\centering
\caption{Cross-partition comparison of the FedProx-vs-FedAvg within-pair difference under each system-heterogeneity condition. The engineered partition contains substantial class skew across clients (every minority class held by two clients); the IID partition contains no inter-client class skew. The proximal anchor's mechanism predicts a positive $\Delta$ on the engineered partition and $\Delta \approx 0$ on IID.}
\label{tab:system-het-cross-partition}
\small
\begin{tabular}{lcccc}
\toprule
& \multicolumn{2}{c}{Engineered partition} & \multicolumn{2}{c}{IID partition} \\
\cmidrule(lr){2-3} \cmidrule(lr){4-5}
Condition & $\Delta$ & Wins & $\Delta$ & Wins \\
\midrule
C0 (no system het) & $+0.0069$ & 7/10 & $-0.0061$ & 5/10 \\
C1 (fixed stragglers) & $+0.0133$ & 8/10 & $+0.0102$ & 6/10 \\
C2 (random stragglers) & $+0.0174$ & 7/10 & $-0.0131$ & 3/10 \\
\bottomrule
\end{tabular}
\end{table}


Figure~\ref{fig:system-het-convergence} shows the trajectories
side-by-side: the top row (engineered) consistently displays a
visible FedProx-FedAvg gap, while the bottom row (IID) shows the
two algorithms' trajectories essentially superimposed. The visual
contrast between the two rows is the most direct illustration of
the mechanism's data-dependence: FedProx helps where there is class
skew, and is inert where there is not.

\subsection{Partial client participation}
\label{sec:partial-participation}

The conditions reported above vary the second canonical FL parameter
of system heterogeneity --- the local-epoch budget $E$ --- while
holding the per-round participation fraction fixed at $C = 1.0$ (every
client participates every round). The third canonical FL parameter,
$C$, is itself an axis of system heterogeneity widely tested in the
literature (\citealt{mcmahan2017fedavg} §3.1 varies
$C \in \{0.1, 0.2, 0.5, 1.0\}$; \citealt{li2020fedprox} §5.3 tests
FedProx under partial participation on Synthetic and FEMNIST;
\citealt{kairouz2021advances} lists $C$ as a defining FL setup
parameter). In real federated deployments, not all clients respond
every round, so the full-participation setting reported above is the
cleanest baseline but the least operationally realistic case. We
additionally report a partial-participation condition at $C = 0.5$
(four of the seven clients sampled uniformly at random per round,
computed as $\max(1, \lfloor 0.5 \times 7 \rfloor) = 4$), holding all
other hyperparameters identical to the engineered C0 baseline.

Table~\ref{tab:partial-participation} reports the result. Under
partial participation the within-pair $\Delta$ is
\TODOhpc{[fill in: $X.XXX$ macro-F1, $Y/10$ paired wins once HPC
results land]}, [direction: descriptively positive / comparable to /
descriptively negative versus] the $C = 1.0$ baseline of $+0.007$
macro-F1. The partial-participation setting introduces additional
aggregation-pool variance per round (fewer clients contributing,
varying subsets across rounds), which on mechanism grounds is
expected to interact with the proximal anchor's drift-control
mechanism in a regime not directly tested by the local-epoch
heterogeneity conditions above. The descriptive result here covers
the second canonical FL system-heterogeneity axis and complements the
local-epoch heterogeneity reported in
Sections~\ref{sec:system-het-engineered}--\ref{sec:system-het-iid}.

\input{T14_partial_participation}

\subsection{Replication of the canonical Li et al.\ §5.2 protocol}
\label{sec:asymmetric}

The H1/H2 tests above use a symmetric inclusion policy in which both
FedAvg and FedProx aggregate over all clients including stragglers.
The canonical asymmetric protocol of \citet{li2020fedprox} §5.2, in
contrast, has FedAvg drop straggler updates entirely while FedProx
includes them as $\gamma$-inexact partial work. The original paper
reports a substantial FedProx advantage under this asymmetric setup.

To verify that this canonical setup reproduces qualitatively on
DermaMNIST, we additionally ran the FedAvg-drop arm on the engineered
partition with random stragglers (the C2 schedule). Combined with the
existing FedProx-include arm from the symmetric C2 data on the same
ten seeds, this gives the within-pair comparison in
Table~\ref{tab:asymmetric-replication}.

% T09 -- Li 2020 §5.2 asymmetric protocol replication (auto-generated)
\begin{table}[!htbp]
\centering
\caption{Replication of the asymmetric straggler protocol of \citet{li2020fedprox} §5.2, in which FedAvg drops straggler updates (its canonical behaviour) while FedProx includes them as $\gamma$-inexact partial work. Values are mean $\pm$ standard deviation of test macro-F1 across $n = 10$ paired seeds on the engineered partition with random stragglers (4 of 7 clients per round). $\Delta_{\mathrm{total}}$ is the within-pair FedProx-include minus FedAvg-drop difference; $p$ is the two-sided paired Wilcoxon $p$-value. \textbf{Bold} marks contrasts significant at $\alpha = 0.05$.}
\label{tab:asymmetric-replication}
\small
\begin{tabular}{lcccc}
\toprule
Algorithm (aggregation policy) & Test macro-F1 & $\Delta_{\mathrm{total}}$ & $p$ (Wilcoxon) & Wins \\
\midrule
FedAvg (drops stragglers) & $0.4423 \pm 0.0192$ & --- & --- & --- \\
FedProx (includes stragglers) & \textbf{$0.5102 \pm 0.0168$} & \textbf{$+0.0679 \pm 0.0175$} & \textbf{$0.002$} & 10/10 \\
\bottomrule
\end{tabular}
\end{table}


The FedProx-include configuration outperforms FedAvg-drop by a mean
within-pair difference of $+0.068$ macro-F1, with FedProx winning
all 10 paired seeds. This reproduces the qualitative finding of
\citet{li2020fedprox} §5.2 on the DermaMNIST task: under the
canonical asymmetric straggler protocol, the FedProx configuration
produces a substantially higher macro-F1 than the FedAvg
configuration.

\subsection{Summary}
\label{sec:system-het-summary}

On the engineered partition the within-pair FedProx-vs-FedAvg
difference rises monotonically with straggler severity (C0:
$+0.007$, C1: $+0.013$, C2: $+0.017$), consistent with the canonical
\citet{li2020fedprox} motivation. The C2 per-class breakdown localises
the advantage to the melanoma class ($\Delta = +0.085$ macro-F1, 10/10
paired wins). Under partial client participation ($C = 0.5$,
Section~\ref{sec:partial-participation}) the descriptive result is
\TODOhpc{[fill in once HPC data lands]}. The canonical asymmetric
protocol of \citet{li2020fedprox} §5.2 reproduces qualitatively:
$\Delta_{\text{total}} = +0.068$ macro-F1 with FedProx winning all 10
paired seeds. On the IID partition the within-pair $\Delta$ is small
and not consistently positive under any of the three straggler
conditions, in line with the proximal-anchor mechanism's prediction
that the regulariser is inert when no inter-client class skew exists.

\paragraph{Scope and limitations.}
The within-pair $\Delta$ magnitudes on the engineered partition
($+0.007$--$+0.017$ macro-F1) are smaller than the effect sizes
originally reported in \citet{li2020fedprox} on Synthetic, FEMNIST,
Shakespeare, and Sentiment140 benchmarks. This is consistent with
the small absolute magnitude observed in the statistical-heterogeneity
arm under the same partition and runtime, and is plausibly attributable
to the smaller dataset, more challenging architecture, and
medical-imaging difficulty discussed in
Section~\ref{sec:headline-summary}. The C1 fixed-straggler condition
on the engineered partition is structurally coupled with the clinical
role assignment (C5 holds the melanoma/vascular pair); its descriptive
result should be read with that confound in mind. The asymmetric
protocol of Section~\ref{sec:asymmetric} provides a single point of
comparison against the canonical Li et al.\ setup but is not a primary
analysis; it is included to demonstrate qualitative reproducibility of
their §5.2 result on DermaMNIST. The straggler schedules are
deterministic functions of the seed and have been verified by direct
reconstruction from the JSON provenance --- the empirical per-client
straggler frequency, $E_i$ range, and round counts match the schedule
specifications exactly.
